\chapter{术语表}
\markboth{术语表}{术语表}

\begin{ledetext}
本表收录正文中出现的核心术语，按汉语拼音排序。每条后面标注它在哪一章展开。看到不认识的词回来查一眼，然后跳回正文，看它是怎么长出来的。
\end{ledetext}

\begin{multicols}{2}

\glossletter{A}
\glossitem{72 法则}{}{72 除以年化收益率，约等于本金翻倍所需年数。按 3.5\% 算，需要约 20.6 年。}{03}

\glossletter{B}
\glossitem{保证金追缴}{margin call}{价格下跌触发追加保证金，拿不出就被强平。判断正确但活不到被证明的那天，就是这么来的。}{14}
\glossitem{被动投资}{}{按指数权重机械买入。2025–2026 年全球规模首次超过主动管理，价格由谁来发现成了新问题。}{11}
\glossitem{遍历性}{ergodicity}{集合平均不等于时间平均。一百个人各玩一次和一个人连玩一百次，是两件完全不同的事。}{22}
\glossitem{波动损耗}{volatility decay}{涨 10\% 再跌 9.09\% 回到原点，三倍杠杆产品却亏了 5.5\%。杠杆 ETF 只适合极短期。}{14}

\glossletter{C}
\glossitem{储蓄率}{}{收入减支出的比例。这是全书里唯一一个你能百分之百控制的变量。}{19}
\glossitem{存款保险}{}{中国 50 万元以内本息全额偿付。它真正的作用不是事后赔钱，是事前掐断挤兑的念头。}{02}

\glossletter{D}
\glossitem{贷款创造存款}{}{银行放贷时同时凭空写下一笔资产（贷款）和一笔负债（存款），新货币由此诞生。全书最反常识、也最重要的一条。}{07}
\glossitem{到期收益率}{YTM}{按当前价格买入并持有到期的年化回报。价格与它永远反向变动。}{10}

\glossletter{F}
\glossitem{法定货币}{fiat money}{国家指定、以税收和法偿性为支撑的货币。它的价值不来自黄金储备，来自你必须用它交税。}{01}

\glossletter{H}
\glossitem{货币乘数}{}{过时的教科书模型：认为央行投放准备金后银行按倍数放大。现代银行体系里因果是反过来的。}{07}

\glossletter{J}
\glossitem{挤兑}{bank run}{一个自我实现的协调博弈：我不去取钱的前提是别人也不去。硅谷银行证明了它在社交媒体时代能有多快。}{02}
\glossitem{净息差}{NIM}{银行放贷收益与负债成本之差。中国商业银行 2026 年二季度末为 1.41\%，处于历史低位。}{02}
\glossitem{净值}{net worth}{资产减负债。房子按市价记，房贷按余额记，中间的差额才是你的。收入不是财富，净值才是。}{19}
\glossitem{久期}{duration}{债券对利率的敏感度。久期 8 年，利率升 1\%，价格大约跌 8\%。硅谷银行死在这两个字上。}{10}

\glossletter{L}
\glossitem{量化宽松 / 紧缩}{QE / QT}{央行买入或卖出债券以影响长端利率。QE 灌的是银行间的准备金池，不会自动流进你的口袋。}{07}

\glossletter{M}
\glossitem{名义利率 / 实际利率}{}{实际利率 ≈ 名义利率 − 通胀。存款 2\% 而通胀 3\%，你的钱在变多，购买力在变少。}{03}
\glossitem{明斯基时刻}{}{稳定本身孕育不稳定。日子越安稳，人越敢借钱，系统越脆弱。}{17}

\glossletter{Q}
\glossitem{期权}{option}{买方付权利金换取权利：损失有限，收益可观。卖方相反，像在压路机前捡硬币。}{15}
\glossitem{期限转换}{maturity transformation}{银行的核心魔法：把随时可取的存款变成三十年才还完的房贷。它创造价值，也制造脆弱。}{02}
\glossitem{前瞻指引}{forward guidance}{央行用语言影响预期。2026 年美联储主席沃什明确表示要淡化它，这是沟通范式的转向。}{08}
\glossitem{权益乘数}{}{总资产 ÷ 净资产。银行通常十几倍，所以 0.8\% 的资产收益率能做出两位数的股东回报。}{14}

\glossletter{R}
\glossitem{人力资本}{human capital}{你未来收入的贴现值。对大多数人来说这是最大的一笔资产，而它不在任何一张对账单上。}{19}

\glossletter{S}
\glossitem{三元悖论}{}{固定汇率、资本自由流动、独立货币政策，三者最多得其二。香港、欧元区、中国各自选了不同的两个。}{12}
\glossitem{剩余索取权}{residual claim}{股东的法律身份：付完供应商、工资、利息、税，剩下的才归你。排在最后，所以风险最高、长期回报也最高。}{11}
\glossitem{市盈率}{PE}{股价除以每股收益。它不是独立真理，只是贴现公式的一个简写。}{11}
\glossitem{收益率曲线}{}{不同期限国债收益率连成的线。倒挂曾是可靠的衰退先行指标，2022–2023 年它失灵了，争议至今未平。}{09}
\glossitem{私募信贷}{private credit}{2025 年全球规模约 2.3 万亿美元，利率 8\%–12\%。中金称之为“2 万亿美元的灰犀牛”。}{16}

\glossletter{T}
\glossitem{套息交易}{carry trade}{借低息货币买高息资产。利差是慢慢赚的，汇率是一夜之间还回去的。}{12}
\glossitem{特里芬难题}{}{世界要用美元做储备，美国就得持续输出美元；而长期逆差又会侵蚀对美元的信心。}{12}

\glossletter{X}
\glossitem{系统性风险}{}{分散不掉的那部分。2008 年的教训是：危机中所有相关性都趋近于 1。}{13}
\glossitem{现值与贴现}{present value}{未来的钱换算成今天值多少。贴现率一变，全世界所有资产的价格都要重算。}{03}
\glossitem{信用利差}{credit spread}{同期限下信用债与国债的收益率之差，用来补偿违约风险。它比股市更早知道事情不对。}{13}
\glossitem{信用违约互换}{CDS}{给别人的债买的保险，而且你可以给不属于你的房子买火险。2008 年 AIG 就栽在这上面。}{15}
\glossitem{序列风险}{sequence risk}{同样的平均收益，顺序不同结果天差地别。在退休提取阶段，头几年的大跌是致命的。}{22}

\glossletter{Y}
\glossitem{一级 / 二级市场}{}{公司发新股拿到钱的地方 / 投资者之间转手的地方。你在二级市场买股票，公司一分钱也拿不到。}{11}
\glossitem{影子银行}{}{干着银行的三件事，却不在银行监管框架内，也没有存款保险。它不违法，它是金融体系的正常器官。}{16}
\glossitem{预期损失}{}{违约概率 × 违约损失率 × 敞口。小微企业融资贵不是银行黑心，先是数学，然后才是摩擦。}{13}

\glossletter{Z}
\glossitem{再平衡}{}{按纪律把偏离目标的仓位调回来。它强迫你在下跌时买入、上涨时卖出，而这两件事靠意志力做不到。}{21}
\glossitem{政策利率}{}{央行直接设定的那个价格。中国以 7 天期逆回购利率（1.4\%）为主，美国是联邦基金利率目标区间（3.50\%–3.75\%）。}{08}
\glossitem{滞胀}{}{增长停滞与通胀并存。2026 年的组合是供给冲击加高利率加高债务，这是它危险的原因。}{18}
\glossitem{准备金}{reserves}{商业银行存在央行的钱，只在银行之间清算时使用。QE 创造的是它，不是你账户里的存款。}{07}
\glossitem{资产负债表衰退}{}{企业和居民的目标从赚钱变成还债，利率再低也不借。日本用三十年演示了一遍。}{09}
\glossitem{最后贷款人}{}{白芝浩原则：危机时对有偿付能力的机构、以惩罚性利率、凭优质抵押品无限放款。}{08}
\glossitem{LPR}{}{贷款市场报价利率。18 家报价行在政策利率基础上加点报出，你的房贷 = LPR ± 合同里固定的加点。}{09}
\glossitem{M0 / M1 / M2}{}{流通中的现金 / 加上企业活期存款 / 再加上定期与储蓄存款。中国 M2 三百多万亿元，M0 只有十几万亿——95\% 以上的钱从没当过纸。}{01}
\end{multicols}
\cleardoublepage
